\documentclass[11pt,a4paper]{article}
\usepackage[T1]{fontenc}
\usepackage{lmodern,amsmath,amssymb,amsthm,mathtools,booktabs,microtype}
\usepackage[margin=27mm]{geometry}
\usepackage[colorlinks=true,linkcolor=blue,citecolor=blue,urlcolor=blue]{hyperref}
\setlength{\emergencystretch}{2em}
\newtheorem{theorem}{Theorem}[section]
\newtheorem{lemma}[theorem]{Lemma}
\newtheorem{proposition}[theorem]{Proposition}
\newtheorem{corollary}[theorem]{Corollary}
\newtheorem{definition}[theorem]{Definition}
\newcommand{\Q}{\mathbb Q}\newcommand{\Z}{\mathbb Z}
\DeclareMathOperator{\Tr}{Tr}\DeclareMathOperator{\Nm}{N}
\title{A rational conic criterion using the actual detector trace}
\author{}\date{17 November 2026}
\begin{document}\maketitle
\begin{abstract}
For normalized elliptic detectors, one rational square depending only
on the actual reciprocal trace excludes every relevant negative odd
pole. Its radicand is positive at every prepared character and equals
four at the reference product. The auxiliary curve in trace and root
is a rational conic. The test is locally square at every fixed
preparation place and selected contact of the existing completeness
construction. Simultaneous rational realization at the two actual
detector traces remains an additional arithmetic requirement; if it
holds, the finite-algebra constructions have twenty-one and nine
total variables. No such realization or new defining count is proved.
\end{abstract}

\section{The trace square}

Let \(\mathcal S\) be the fixed exceptional set of the source, with
\(2\in\mathcal S\), and let
\[
 w\in\Q^\times,\qquad V\in\Q^\times,\qquad b=\frac{4w}{V^2}.
\]
Assume that every odd valuation of \(w\) outside \(\mathcal S\)
is positive. In the elliptic application \(V=\Tr_{K/\Q}(W^{-1})\),
where \(K=\Q(\sqrt2)\). The normalized general weight
\(w=u/(1+u)^2\) and the conditional weight
\(F(t)/(4(t-1)^8)\) have this property on their retained guarded
domains \cite{Native,Normalized}.

Define
\begin{equation}\label{eq:conic}
                       \mathcal C=V^2+w(2-V).
\end{equation}

\begin{theorem}\label{thm:pole}
If \(\mathcal C\) is a rational square, then \(b\) has no negative
odd valuation outside \(\mathcal S\). This implication does not
require any residue condition on the individual reciprocal slots.
At every prepared character where \(w>0\) and the two real slots of
\(W\) are at least one, \(\mathcal C\) is strictly positive.
At the algebraic reference product \(W=1\), its value is exactly four.
\end{theorem}
\begin{proof}
Suppose \(v_p(b)<0\) is odd, with \(p\notin\mathcal S\).
Then \(\nu=v_p(w)\) is odd, hence positive by hypothesis, and
\(2v_p(V)>\nu\). In particular \(v_p(V)>0\), and because
\(p\ne2\), the number \(2-V\) is a unit. The two summands in
\eqref{eq:conic} have distinct valuations \(2v_p(V)>\nu\) and
\(\nu\). Thus \(v_p(\mathcal C)=\nu\) is odd, contradicting
rational squareness. This argument also excludes the possibility
\(\mathcal C=0\) at such a bad pole.

At a prepared real character, the two reciprocal slots are in
\((0,1]\), so \(0<V\le2\). Both terms in \eqref{eq:conic} are
nonnegative and the first is positive. If \(V=2\), including the
reference product, the value is four. No limiting or nonzero
assumption on \(2-V\) is needed.
\end{proof}

The conclusion excludes exactly the primes used in the retained source
height argument. It is not an all-prime statement at two. The original
guards \(w\ne0\) and \(V\ne0\) are retained.

Writing \(W=A+B\sqrt2\) and \(N=A^2-2B^2\), one has
\begin{equation}\label{eq:cleared}
 \mathcal C=\frac{4A^2+2wN(N-A)}{N^2}.
\end{equation}
Thus the new test introduces no denominator divisor beyond those
already used for the weight and the detector norm. In particular, no
norm-of-reciprocal value is an additional prescribed arithmetic datum.

\section{Local preparation with the actual trace}

\begin{proposition}\label{prop:local}
The existing normalized general and conditional completeness
preparations give the following local square conditions for
\eqref{eq:conic}:
\begin{enumerate}
\item At each selected contact prime, where \(v_p(w)>0\) and
\(V\equiv1\) or \(2\pmod p\), the value \(\mathcal C\) is
a nonzero square in \(\Q_p\).
\item At every fixed preparation place, the source may arrange that
\(\mathcal C\) is a nonzero local square by taking \(V\) sufficiently
close to two, with the nonzero physical weight in its prescribed
neighborhood.
\end{enumerate}
Both assertions retain the actual detector trace, point tuple, common
slope and comparison scales.
\end{proposition}
\begin{proof}
At an odd selected contact, \(V\) is a unit and
\[
             \frac{\mathcal C}{V^2}
             =1+\frac{w(2-V)}{V^2}\in1+p\Z_p.
\]
Hensel's lemma gives a square. In particular its residue is one or
four according to the two possible trace residues.

At a fixed odd place, make the valuation of the second term in the
displayed ratio positive. At two, make it at least three. These
conditions are possible by taking \(V\) sufficiently close to two;
the weight stays in a fixed neighborhood of a nonzero value, or in
the stated conditional preparation tends to zero. The familiar local
square neighborhoods \(1+p\Z_p\), for odd \(p\), and
\(1+8\Z_2\) prove the assertion. The general source prepares its
scalar functions first and then the slope, so it allows this arbitrary
fixed precision without changing the order of choices.
\end{proof}

The proposition does not settle local square conditions at new primes
dividing detector numerators with even physical valuation. It also
does not imply rational squareness or simultaneous global realization.

\section{Rational parametrizations}

\begin{proposition}\label{prop:param}
For fixed \(w\ne0\), the conic \(Y^2=V^2+w(2-V)\) has the
rational point \((V,Y)=(2,2)\). On the chart \(1+2k\ne0\) its
rational parametrization is
\begin{equation}\label{eq:param}
 V=\frac{2-k^2w}{1+2k},\qquad Y=V+kw.
\end{equation}
Every rational point with \(1+2(Y-V)/w\ne0\) lies on this chart.
The exceptional value \(k=-1/2\) can occur only when \(w=8\);
then the conic splits as \(Y^2=(V-4)^2\).
An alternative chart from the line through \((2,2)\) is
\[
 V=\frac{2(t-1)^2-w}{t^2-1},\qquad Y=2+t(V-2),
 \qquad t^2\ne1.
\]
\end{proposition}
\begin{proof}
Substitute \(Y=V+kw\), cancel the nonzero \(w\), and obtain
\((1+2k)V=2-k^2w\). The exceptional case and converse follow
directly. Substitution of \(Y=2+t(V-2)\), followed by removal of
the reference-point factor \(V-2\), gives the second chart.
\end{proof}

These are arithmetic parametrizations, not new unrestricted coordinates
in the finite-algebra specification. A solution must match the actual
\(V\) and \(w\) coming from the retained point tuple and common
slope. The reference value \(W=1\) need not itself be an actual source
tuple.

The simple slice \(Y=2-V\) is equivalent, away from \(V=2\), to
\[
                         w=\frac{4(1-V)}{2-V}.
\]
On the prepared real domain \(w>0\), \(0<V\le2\), this requires
\(0<V<1\). It is therefore outside the usual completeness
neighborhood \(V>1\), although it may be investigated using a
different real preparation. Local compatibility at a contact with
\(V\equiv1\) alone does not provide that new real preparation.

\section{Whole-character encoding and exact coordinate count}

\begin{definition}
The simultaneous trace-conic property for a retained source means
that every required integer input has a prepared actual rational
point tuple and one common slope, with its original scales and all
point, difference, scalar, weight and trace guards, such that both
values
\[
                     V_j^2+w_j(2-V_j),\qquad j=0,1,
\]
are rational squares. The property includes the whole-character real
and guard locus of the source. In a conditional source the required
inputs are the nonzero integers.
\end{definition}

\begin{theorem}\label{thm:census}
Under the source's retained arithmetic premises and its simultaneous
trace-conic property, the general source gives a defining polynomial
with twenty-one total variables, and either conditional source gives
one with nine total variables. Without that property, the construction
proves soundness only. The established complete defining counts are
unchanged.
\end{theorem}
\begin{proof}
Let \(p=9\) in the general source and \(p=3\) in a conditional
source. Retain the \(p\) rational point pairs and their free point
algebra of rank \(4^p\). Replace each two-coordinate pole block
by the finite quadratic test \eqref{eq:conic}. If its rational
expression in the localized point algebra is \(N_j/D_j\), adjoin
the monic generator
\[
                              Y_j^2=N_jD_j.
\]
The original guard product includes all syntactic denominators.
At every prepared character Theorem~\ref{thm:pole} makes the radicand
positive and nonzero. Include the numerator and denominator norms in
the scalar guard product, as required by the existing encoder; these
are automatic nonzero guards on the prepared locus and require no
new rational coordinate. The two extensions are finite etale and
totally real there. Their joint rank is
\[
                          4^p\cdot2^2=4^{p+1}.
\]
The one-coordinate separator and determinant descent of \cite{Unit}
therefore apply, with no additional point, square, height-recovery or
nonzero-witness coordinate.

Every rational solution excludes the same negative odd poles by
Theorem~\ref{thm:pole}. The unchanged source extraction and height
argument give soundness, including the same exceptional set and fixed
defining-parameter qualifications. The simultaneous trace-conic
property extends the selected rational character to both quadratic
generators and proves completeness. The total number of unrestricted
coordinates is
\[
              1\text{ input}+2p\text{ point coordinates}
                +1\text{ slope}+1\text{ separator}=2p+3.
\]
The conditional source's original final input factor includes zero.
\end{proof}

All comparisons in the general source, including its Lucas relation
and scalar-height argument, are retained. No ordinary or weak abc
hypothesis is introduced into the general criterion. No degree or
coefficient bound for the complete resulting determinant is asserted.

\section{Trace-dependent classes with no lost contact primes}

\begin{proposition}\label{prop:class}
In Theorem~\ref{thm:pole}, replace rational squareness of
\(\mathcal C\) by that of \(\delta\mathcal C\), where \(\delta\)
is a retained rational source function whose valuation is even at every
hypothetical negative odd pole of \(b\). The pole implication is
unchanged. If \(\delta\) is positive at every prepared character
and all its denominators are retained guards, the finite-algebra
argument has the same coordinate count.

In particular, both functions
\begin{equation}\label{eq:traceclass}
       \delta_{\mathrm{lin}}(V)=1+24V,\qquad
       \delta_{\mathrm{quad}}(V)=1+3(V-1)^2
\end{equation}
are admissible for the general and conditional sources, without any
height-support budget. They introduce no lost contact prime. Each is
locally square at every fixed preparation place and every selected
contact of the retained constructions, without a splitting assumption.
\end{proposition}
\begin{proof}
At a bad pole the valuation of \(\mathcal C\) is odd, and adding
the even valuation of \(\delta\) cannot change that parity.
For \eqref{eq:traceclass}, one has \(v_p(V)>0\) at a bad pole,
so the linear and quadratic functions reduce to one and four,
respectively. Both are strictly positive on the prepared domain
\(0<V\le2\), and their denominators are powers of that of \(V\),
already retained. At the reference trace their values are forty-nine
and four. Arbitrary fixed local precision near \(V=2\) makes each
a local square, including at two and seven.
At a selected contact the linear function reduces to twenty-five or
forty-nine, and the quadratic function to one or four. The selected
primes are greater than \(76176\), so these are nonzero square
residues. No splitting assumption is needed. No contact is lost in the
soundness proof, and the original height constants and exponents
are unchanged.
\end{proof}

More generally, the same assertions hold for
\[
          \delta_r(V)=r^2+(4-r^2)(V-1)^2,
\]
where \(r\) is a fixed positive \(\mathcal S\)-unit with \(r<2\).
Its values at traces zero and two are four, its value at trace one is
\(r^2\), and it is strictly positive on the entire real line.
The choice \(r=1\) has integral coefficients. The product
\(\delta_{\mathrm{quad}}(V)\mathcal C\) has total degree four in
\((V,w)\). The linear class gives the smaller cubic radicand
\begin{equation}\label{eq:cubic}
 (1+24V)\mathcal C
 =24V^3-24V^2w+V^2+47Vw+2w,
\end{equation}
with five monomials and coefficient height forty-seven. These are
degrees and coefficients of the displayed trace radicands, not bounds
for the complete defining determinant.

For fixed \(w\), the cubic curve has the rational point
\((V,Y)=(2,14)\). The change \(X_c=1+24V\), \(Y_c=24Y\)
puts it in the form
\[
 Y_c^2=X_c^3-(2+24w)X_c^2+(1+1176w)X_c,
\]
where the reference point is \((49,336)\) and \((0,0)\) is a
rational point of order two when the curve is nonsingular.
The discriminant is
\[
                    9216w(w-8)(1+1176w)^2.
\]
Thus every general prepared weight \(0<w\le1/4\) gives a smooth
elliptic curve. This model retains a prescribed abscissa
\(X_c=1+24V\) when used in the actual detector problem.

In particular, the finite menu
\(\{1,1+24V,1+3(V-1)^2\}\) is available separately in each block.
One may use any fixed finite menu of such positive admissible
functions, by a finite product of branch encoder polynomials sharing
the same unrestricted coordinates. This increases polynomial size,
but introduces no additional rational variable. It provides further
square classes for a future actual-trace realization; it is not itself
such a realization.

\begin{thebibliography}{9}\raggedright
\bibitem{Native}\emph{A guarded rational-square criterion for normalized
elliptic detectors},
\nolinkurl{papers/pole_parameters/native_square_criterion_20261116_v53.tex}.
\bibitem{Normalized}\emph{Quadratic reweightings of native elliptic
detectors},
\nolinkurl{papers/pole_parameters/reweighted_elliptic_detectors_20261116_v53.tex}.
\bibitem{Unit}\emph{A single norm determinant for the 25- and 13-variable
constructions},
\nolinkurl{papers/finite_algebra/unit_norm_20261114_v28.tex}.
\bibitem{Scalar}\emph{Intrinsic congruent-curve poles from a fixed scalar
rescaling},
\nolinkurl{burungale_scalar_transfer_20261117_v56.tex}.
\end{thebibliography}
\end{document}
